\section{Hybrid Retrieval Strategies}
\label{sec:gr-hybrid}

Graph and vector retrieval have complementary strengths, so effective systems
combine them rather than choosing one: dense similarity is used to enter the graph
and locate relevant entities, while graph traversal expands that entry point into
a connected, relationally coherent context~\cite{peng2024graphrag}. LightRAG
embodies this by integrating graph structures with vector representations in a
single index, which it reports improves both retrieval accuracy and response time
over conventional RAG~\cite{guo2024lightrag}.

Hierarchical summarization can also be seen as a form of hybrid retrieval over
levels of abstraction rather than over data models. RAPTOR, for example,
recursively embeds, clusters, and summarizes text to build a tree, and retrieving
from different levels of that tree integrates information at multiple granularities,
yielding significant gains over flat retrieval~\cite{sarthi2024raptor}. Whether the
structure is a knowledge graph or a summary tree, the common principle is that
combining semantic similarity with explicit structure retrieves more complete and
better-connected evidence than similarity alone~\cite{peng2024graphrag,
sarthi2024raptor}.
